\section{Validation Against Runtime Artifacts}
\label{sec:validation}

We evaluate the contracts against one public trace corpus and one public
post-oracle conformance fork, and retain one local instrumented study as
illustrative evidence.  The public, sealed Synthetic Players capsule contains
an ActiveGraph SQLite store.  The public activegraph-bridge fixture executes an
actual bridge fork after a committed recorded fixture-oracle effect.  The local
executive study records explicit boundary verification but has no public source
remote.  All three contain synthetic experimental content; ``real'' here means
actual runtime artifacts rather than invented shape fixtures.  None is
independent of the ActiveGraph family, and we do not claim human behavioral
validity or real-model performance.

\subsection{Synthetic Players capsule}

The public source is pinned to revision \code{82556ef9dae3}.  The compressed
engine store begins SHA-256 \code{7b7fdfe64b5b}, and the capsule checksum
ledger begins SHA-256 \code{6d80f3b19b92}.  The evidence ledger records the
complete revision and digests.
Running the public verifier reports:

\begin{quote}
CAPSULE VERIFICATION PASS --- 4,919 archived Phase 3--5 runs verified
(4,916 confirmatory + 3 legacy diagnostics).
\end{quote}

The decompressed database contains 5,540 runs and 311,756 events.  Event types
include 36,251 \code{llm.requested}, 36,227 \code{llm.responded}, and 30,397
\code{decision.parsed} records.  The difference between request and response
counts is reflected rather than repaired by the adapter.

The two run counts have different denominators.  The store contains 5,505
completed and 35 incomplete runs.  The public verifier covers 4,916
confirmatory Phase 3--5 runs plus three completed legacy diagnostics.  The
remaining 586 completed runs are outside that replay contract; the 35
incomplete runs are also outside it.  Thus $4{,}919+586+35=5{,}540$.

\subsubsection{Replay grades}

All 5,540 logs grade Observed, and zero are verified \emph{from the run log
alone}.  This does not contradict the passing capsule verifier.  The capsule
attestation is external to each run's event stream.  A grading function that
claims to depend on the log alone cannot silently import that attestation.
This is a useful negative result: sealed-artifact reproducibility and
self-describing run replayability are distinct properties.

\subsubsection{All-cut results}

Table~\ref{tab:syntheticcuts} reports every cut.  The projection total is
$311{,}756 + 5{,}540 = 317{,}296$, because each run of length $n$ has $n+1$
cuts.

\begin{table}[ht]
\centering
\small
\begin{tabular}{@{}lrrr@{}}
\toprule
Property & Sound & Conditional & Unsound \\
\midrule
ProjectionReplay & 317,296 & 0 & 0 \\
ExternalContinuation & 160,076 & 120,969 & 36,251 \\
CounterfactualWorld & 150,230 & 24,172 & 142,894 \\
\bottomrule
\end{tabular}
\caption{All-cut verdicts for 5,540 Synthetic Players logs.}
\label{tab:syntheticcuts}
\end{table}

Every projection cut is sound under the specified domain-projection
observation.  The 36,251 ExternalContinuation-unsound cuts equal the number of
open/cut-through LLM requests.  Direct sequence audit explains the equality:
36,227 requests are immediately followed by their linked response, so each
creates exactly one cut between request and response; the remaining 24 requests
are the final event of failed runs, so each also creates exactly one open cut.
Conditional continuation cuts retain a committed OneShot oracle interaction
and require serving the recorded result without another call.
CounterfactualWorld is stricter: 4,923 runs have at least one Unsound cut
because a discarded oracle interaction already occurred or a boundary is
ambiguous.

The 160,076 Sound ExternalContinuation cuts validate only the
no-inherited-effect case.  A direct per-run ordinal query confirms that every
one lies before the first classified external request, including all cuts in
the 617 runs with no classified request.  The Synthetic Players corpus alone
therefore does not empirically validate a post-oracle continuation,
target-environment attestation, or zero-reexecution receipt.  Generated
properties and named fixtures exercise post-oracle verdicts and typed
no-reexecution obligations; the public conformance fork below separately
exercises their runtime discharge under a fixture trust root.

The validator leaves domain and infrastructure types unclassified, including
\code{behavior.started} and \code{decision.parsed}.  It likewise leaves
\code{infra.*}, \code{round.requested}, and \code{run.completed} unclassified.  In particular,
\code{round.requested} is not misclassified as an external request merely
because its name ends in \code{requested}.  Continuation verdicts are therefore
relative to the audited profile declaration that these reported types are
domain or infrastructure signals rather than hidden external effects.

\subsubsection{Observed forks}

The SQLite run table contains 121 child runs from 41 distinct parents and no
nested child in this sample.  All forks cut at a \code{round.played} event.  The
audit compares 16,625 parent-prefix events with the beginning of their child
traces and finds zero mismatches in event ID, type, actor, payload, frame, cause,
or timestamp.  Every cut event resolves.  Run identity and \code{events.seq}
are excluded because they are expected to differ: in this store schema,
\code{seq} is the database-wide autoincrement key used for ordering, not a
run-local field.  The per-run normalized sequence is derived after export and
was therefore not an additional physical field available for comparison.

None of the 121 retained prefixes contains a classified external request, and
none contains an unresolved request.  The observed forks therefore satisfy the
current ExternalContinuation precondition for a domain-only cut family.  This
is strong evidence for cheap structural branching where the runtime actually
forked.  It is not evidence for an effect-boundary fork: the observed sample
does not contain one.

\subsection{Public post-oracle conformance fork}

The companion activegraph-bridge artifact is public at revision
\code{8855d3a9e779}.  The complete
\href{https://github.com/yoheinakajima/activegraph-bridge/tree/8855d3a9e779362f713b08bceb58d7d5db671c7d/evidence/post-oracle-fork-v1}{public post-oracle fixture and verifier}
are revision-pinned.
It records an actual bridge child fork after a committed effect classified as
OneShot footprint plus Recorded replay source.  The source run makes one
deterministic offline fixture-oracle call.  Verification and the child serve
the committed outcome from the record, and the child diverges only after a
changed tool result.

The receipt binds the parent and child run identities, cut boundary, parent log
hash, exact retained-prefix hash, child log hash before receipt emission,
inherited request and outcome identities, response hash, source and target
runtime fingerprints, and target-environment claims.  Its verifier checks those
hashes, the zero-call counter, and an HMAC-SHA256 signature under a configured
trust root whose claims include the child, cut, prefix, and target fingerprint.
The published command reports:

\begin{quote}
POST-ORACLE FORK RECEIPT PASS --- committed oracle served from record;
0 inherited external calls.
\end{quote}

The parent JSONL digest begins \code{884e8ae34296}; the child begins
\code{747305a6774a}.  The receipt file begins \code{da5503a80b92}, and its
canonical internal receipt hash begins \code{5d36df7dc177}.  The verified receipt names inherited request
\code{evt\_008}, its recorded outcome \code{evt\_009}, one source oracle call,
and zero inherited external calls in the fork.

This discharges the prior implementation gap for a public conformance case.  An
inherited recorded oracle outcome can cross a fork boundary without a second
call, and the target environment premise can be authenticated against a
configured trust root.  The oracle is deterministic and offline, and the HMAC
key is deliberately published as a fixture key.  The result does not attest a
real provider, production identity, model quality, or production environmental
fidelity.

\subsection{Illustrative local executive study}

The second artifact contains 12 synthetic vendor due-diligence cases evaluated
under two deterministic offline mock strategies, for 24 runs and 2,592 source
events.  The manifest states that no real LLM or external provider was used and
limits the claim to instrumentation feasibility.  It records fresh-factory
reconstruction and enables boundary verification.

The artifact files are hash-bound.  Hash prefixes for the run store, manifest,
metrics, and paired results are \code{7a38a1e8}, \code{da1cbc3b},
\code{9f9242e1}, and \code{f4d95dbb}, respectively; the evidence ledger records
the complete hashes.  Review exposed that the containing migration-lab had no
commit or remote.  We repaired the local half of that defect: source revision
\code{54dfde4d7379} and release-bundle revision \code{201efa7} bind a redacted
manifest, metrics, pairs, report, compressed run store, and verifier.  The
repository still has no configured remote.  This result remains illustrative,
is excluded from the abstract, and is not externally source-reproducible until
those revisions are published.

\subsubsection{Decision and path equivalence}

The paired artifact reports 66.6667\% final-decision agreement and 0\% ordered
tool-path agreement.  Eight of twelve pairs reach the same decision via
different paths, yielding 66.6667\% same-decision/different-path.  This is not a
model-quality result---both strategies are deterministic mocks.  It is an
instrumented witness that decision equivalence and trace equivalence answer
different questions.  A system that reports only agreement can hide total path
divergence; a system that requires exact paths can reject decisions that are
equal under a task-level observation.

\subsubsection{Grades and cuts}

Each run's source events normalize to two additional evidence facts: fresh
reconstruction and mediated boundary.  The 2,592 source events therefore become
2,640 normalized events.  All 24 logs grade Boundary and all 24 are verified
from explicit log facts.  The 48 additional facts create 48 diagnostic
intra-source positions that cannot be selected as cuts in the atomic source
log.  We therefore report source-event-boundary verdicts as the primary bridge
result and retain every-normalized-position counts as a transformation check.

\begin{table}[ht]
\centering
\small
\begin{tabular}{@{}lrrr@{}}
\toprule
Property & Sound & Conditional & Unsound \\
\midrule
ProjectionReplay & 2,616 & 0 & 0 \\
ExternalContinuation & 1,968 & 264 & 384 \\
CounterfactualWorld & 0 & 264 & 2,352 \\
\bottomrule
\end{tabular}
\caption{Source-event-boundary verdicts for 24 bridge logs.}
\label{tab:bridgecuts}
\end{table}

All 2,616 source-boundary ProjectionReplay cuts are Sound; the diagnostic over
every normalized position is 2,664 Sound.  For completeness, the normalized
ExternalContinuation totals are 1,992 Sound, 288 Conditional, and 384 Unsound;
the CounterfactualWorld totals are 0 Sound, 288 Conditional, and 2,376 Unsound.
Every run contains a OneShot
\code{submit\_recommendation} effect.  No cut is unconditionally Sound for
CounterfactualWorld: before the commit, the effect lies in the discarded
suffix; after the commit, continuation is Conditional on not executing it
again.  The result demonstrates why Boundary grade licenses effect-boundary
operations but cannot promise a world in which discarded actions never
happened.  These counts use the adapter-wide $\Omega_{profile}$, under which
\code{submit\_recommendation} is in scope.  The public bridge v0.2 receipt
records per-effect observable tags, but the legacy executive-study adapter does
not expose such tags to this analysis.  A narrower observable set could change
the world verdicts; we do not report that unmeasured result.

\subsection{Validation status}

\paragraph{V1: compute grade distributions.}
Completed.  The public capsule yields Observed for 5,540 logs because
verification is external to the log; the bridge yields Boundary for 24 logs
because the necessary evidence is explicit.  The contrast supports log-derived
grading rather than profile-derived labels.

\paragraph{V2: inspect real forks and their preconditions.}
Completed for the public capsule and the public conformance fixture.  All 121
observed corpus forks preserve the compared prefix fields and retain no
classified external request.  The companion fixture tests one actual bridge
fork after a committed recorded oracle call and verifies zero inherited
external calls plus an authenticated fixture-environment premise.  It is
conformance evidence, not a real-provider or production-attestor result.

\paragraph{V3: test production confluence conditions.}
Partially completed.  Source audit confirms that ActiveGraph imposes FIFO plus
registration order, which can mask non-confluence.  The property suite bypasses
that order.  A complete inventory of production behavior reads, writes,
emissions, triggers, and relation-order dependencies was not available, so we
cannot certify or falsify a general production non-interference contract.

\paragraph{Combined interpretation.}
The evidence supports deterministic projection, structurally correct observed
prefix copying, and one public post-oracle continuation under a verified
conformance environment.  It does not support an unrestricted claim that any
historical cut is safely executable or environmentally counterfactual.  Those
stronger claims require effect-complete evidence and a property-specific
precondition.
